The Standard Model (SM) is an extremely successful theory that has been rigorously tested at the Large Hadron Collider (LHC) and elsewhere.
Nevertheless it is widely expected that the SM is only an effective field theory (EFT), valid up to some cutoff scale $\Lambda$.
The Standard Model Effective Field Theory (SMEFT) generalizes the SM by adding a complete, but not over-complete basis of operators at every mass-dimension $d$ rather than stopping at $d=4$.\footnote{The SMEFT assumes there are no light hidden states such as sterile neutrinos or an axion, and that the Higgs boson form part of an $SU(2)_w$ doublet with hypercharge $\hyp = 1/2$. Other types of EFTs are possible where these assumptions are relaxed, but we do not consider them here.} 

The counting and classification of operators in the SMEFT has a long history.
Starting with dimension-5 there is a single type operator~\cite{Weinberg:1979sa}, $N_{\rm type} = 1$, and it violates lepton number.
At dimension-6, Ref.~\cite{Grzadkowski:2010es} classified the 76 baryon number preserving ($B$) Lagrangian terms; see~\cite{Buchmuller:1985jz} for earlier work.
The eight baryon number violating ($\slashed B$) terms were previously known~\cite{Abbott:1980zj}, yielding a total of $N_{\rm term} = 84$.
In terms of actual operators rather than terms in the Lagrangian, the counts explode when flavor structure is allowed.
For three generations of fermions, $n_g = 3$, there are $N_{\rm op} = 2499$ independent $B$ operators~\cite{Alonso:2013hga} and 546 $\slashed B$ operators~\cite{Alonso:2014zka}.
Hilbert series methods were applied to the SMEFT in Refs.~\cite{Lehman:2015via, Henning:2015daa, Lehman:2015coa, Henning:2015alf}, providing an elegant way to count the number of operators for arbitrary dimension $d$.
Computing tools $\mathtt{Sim2Int}$~\cite{Fonseca:2017lem}, $\mathtt{DEFT}$~\cite{Gripaios:2018zrz}, $\mathtt{BasisGen}$~\cite{Criado:2019ugp}, $\mathtt{ECO}$~\cite{Marinissen:2020jmb}, and $\mathtt{GrIP}$~\cite{Banerjee:2020bym} were subsequently developed, allowing for automated counting of operators.

Beyond counting operators, work has been done on their explicit forms as well.
Refs.~\cite{Lehman:2014jma, Liao:2016hru} classified the 18 dimension-7 operators.
So far only partial sets of dimension-8 operators exist in the literature.
This includes, however, all of the bosonic operators (in a basis where the number of derivatives is minimized)~\cite{Morozov:1985ef, Hays:2018zze, Remmen:2019cyz}.
Our goal in this work is to find a complete set of dimension-8 operators.
A subtlety in constructing the dimension-8 operator basis is that some operators vanish in the absence of flavor structure.
Our basis contains 231 types of operators.
For comparison, with $n_g = 1$ there are 993 operators, while for $n_g = 3$ there are instead 44807 operators~\cite{Henning:2015alf}.

Although the counting and classification of operators is certainly interesting in its own right, there is also a wide range of phenomenological implications of dimension-8 operators as well.
For some phenomena dimension-8 is the lowest dimension where the interactions become possible.
Most famous among these processes is light-by-light scattering.
Another area where dimension-8 effects have been studied is electroweak precision data (EWPD) where contributions to the $U$ parameter first arise at dimension-8~\cite{Grinstein:1991cd}.
Formally the dimension-6 operators are the leading terms in the EFT expansion.
However there are various scenarios in which this is not the case practically speaking.
Perhaps the most obvious among these is when the interference between the dimension-6 amplitude and the SM amplitude is suppressed or even vanishes.
Additionally there could be a difference in the experimental precision of the measurements being considered~\cite{Dawson:2017vgm}.
Finally, we comment on the structure the renormalization group evolution (RGE) equations of the dimension-8 operators. 

The rest of the paper is organized as follows.
Section~\ref{sec:note} lays down the notation and conventions we use, including the semantics of number of operators versus number of types of operators.
We then discuss how we performed the operator classification in Section~\ref{sec:cls_ops} with the results given in Sec.~\ref{sec:results}.
We briefly explore light-by-light scattering, EWPD, as well as models involving scalar $SU(2)_w$ quartets where there is interesting interplay between dimension-6 and dimension-8 effects in Section~\ref{sec:pheno}.
Additionally we comment on the renormalization group evolution (RGE) of the dimension-8 operators in Sec.~\ref{sec:rge} before concluding in Sec.~\ref{sec:con}.
For convenience we provide tables of dimension-6 and -7 operators in Appendix~\ref{sec:app}.